\documentclass[12pt]{article}
\usepackage[T1]{fontenc}
\usepackage[utf8]{inputenc}
\usepackage{lmodern}
\usepackage{textcomp}
\usepackage{enumitem}
\usepackage{geometry}
\geometry{margin=1in}
\begin{document}
\begin{center}
Answer Key - Business Administration - Graduate Level Assessment 19 \
\small (Answers are ordered exactly as in the question paper.)
\end{center}

\section*{Multiple Choice}
\begin{enumerate}[label=\arabic*.]
\item \textbf{Answer:} E
\\ \textit{Options:} A) 27.5\%; B) 21\%; C) 22.5\%; D) 30\%; E) 24.17\%
\\ \textit{Note:} The correct answer is 24.17\% because the percent markup based on the selling price is calculated by taking the difference between the selling price and the cost ($180 - $136.50 = $43.50), dividing that by the selling price ($43.50 / \$180), and then multiplying by 100 to convert it to a percentage, resulting in approximately 24.17\%.
\item \textbf{Answer:} C
\\ \textit{Options:} A) \$64.12; B) \$83.69; C) \$76.41; D) \$68.24; E) \$74.58
\\ \textit{Note:} The correct answer is $76.41 because the property tax is calculated by multiplying the assessed value of the house, $3,250, by the tax rate of 2.351\%, resulting in $3,250 x 0.02351 = $76.41.
\item \textbf{Answer:} A
\\ \textit{Options:} A) May 15; B) April 15; C) July 14; D) April 14; E) April 13
\\ \textit{Note:} The correct answer is May 15 because a 60-day note starting from March 15 adds 60 days to that date, resulting in a maturity date of May 15.
\item \textbf{Answer:} C
\\ \textit{Options:} A) Economic, emerging, and evolutionary; B) Evolutionary, experimental, and economic; C) Ecological, equitable, and economic; D) Equitable, evolutionary, and emerging; E) Evolutionary, equitable, and economic.
\\ \textit{Note:} The correct answer, "Ecological, equitable, and economic," emphasizes the importance of sustainable practices that balance environmental health, social equity, and economic growth, thereby promoting long-term development over short-term gains.
\item \textbf{Answer:} E
\\ \textit{Options:} A) April 15; B) April 13; C) April 18; D) May 15; E) April 16
\\ \textit{Note:} The loan, made on December 17, is due 120 days later, which calculates to April 16, considering the non-leap year days in each month.
\end{enumerate}

\section*{True / False}
\begin{enumerate}[label=\arabic*.]
\setcounter{enumi}{5}
\item \textbf{Answer:} False
\\ \textit{Options:} A) True; B) False
\\ \textit{Note:} A durable product is typically a high-involvement purchase that requires significant consumer research due to its longevity, higher cost, and impact on the buyer's lifestyle.
\item \textbf{Answer:} False
\\ \textit{Options:} A) True; B) False
\\ \textit{Note:} Jack's commission of 5\% on $1,200 results in earnings of $60 for the week, not $120, as 5% of $1,200 is calculated as $1,200 multiplied by 0.05, which equals $60.
\item \textbf{Answer:} False
\\ \textit{Options:} A) True; B) False
\\ \textit{Note:} The proceeds from a discounted note are calculated by deducting the discount amount from the face value, and with a 3-month note of $850 at a 6% discount rate, the proceeds would be less than $850, not \$860.
\item \textbf{Answer:} True
\\ \textit{Options:} A) True; B) False
\\ \textit{Note:} Role culture is true because it emphasizes a structured environment where specific roles and responsibilities are established, leading to clear expectations and efficient operations within the organization.
\item \textbf{Answer:} False
\\ \textit{Options:} A) True; B) False
\\ \textit{Note:} The statement is false because services are characterized by intangibility, perishability, variability, inseparability, and a lack of ownership, while tangibility is a defining feature of physical goods rather than services.
\end{enumerate}

\section*{Long Form Response}
\begin{enumerate}[label=\arabic*.]
\setcounter{enumi}{10}
\item \textbf{Answer:} Social consensus, temporal immediacy, and proximity are critical factors that shape moral intensity in business decision-making. Social consensus refers to the degree of agreement among stakeholders about what is considered morally acceptable; for instance, a company may hesitate to pursue questionable labor practices if there is a strong community belief against such actions, thereby increasing moral intensity. Temporal immediacy relates to the time frame in which consequences will occur, as decisions with immediate repercussions, such as a product recall, often elicit a stronger moral response than those with delayed effects. Proximity encompasses the closeness of decision-makers to the stakeholders affected by their choices; for example, a manager may feel a heightened moral obligation to act ethically when the consequences directly impact their employees. Together, these factors influence how business leaders evaluate the ethical implications of their decisions, often resulting in a more profound commitment to ethical practices when consensus is strong, consequences are immediate, and proximity to affected parties is high.
\\ \textit{Note:} correct because it effectively identifies and elaborates on the key factors-social consensus, temporal immediacy, and proximity-that influence moral intensity in business decisions, illustrating their impact with relevant examples that demonstrate their significance in shaping ethical considerations.
\item \textbf{Answer:} To analyze the effectiveness of the discount series 30-10-2(1/2)\% compared to 25-15-2\%, we first need to calculate the net price after applying each series of discounts to a hypothetical item priced at $100. For the series 30-10-2(1/2)%, the calculations are as follows: applying 30% first reduces the price to $70, then 10\% off $70 is $63, followed by 2\% off $63 which results in $61.74, and finally applying the additional 0.5\% gives a final price of $61.13. In contrast, the series 25-15-2% starts with a 25% discount, bringing it to $75, then a 15\% discount reduces it to $63.75, and finally a 2% reduction gives a final price of $62.25. Therefore, the series 30-10-2(1/2)\% results in a lower final price of $61.13 compared to $62.25 from the 25-15-2\% series, demonstrating that the former offers better value. Thus, from a financial perspective, the discount series 30-10-2(1/2)\% is more effective than the 25-15-2\% series.
\\ \textit{Note:} correct because it accurately calculates the final net prices of both discount series on a hypothetical \$100 item, allowing for a direct comparison of their effectiveness and value through a clear step-by-step breakdown of the discount application process.
\end{enumerate}

\vfill
\noindent\textit{End of Answer Key}
\end{document}